\chapter*{Introduction}
\addcontentsline{toc}{chapter}{Introduction} 

\openepigraph{The reasonable man adapts himself to the world: the unreasonable one persists in trying to adapt the world to himself. Therefore all progress depends on the unreasonable man.}{George Bernard Shaw, Man and Superman}


\newthought{This} wasn't the book I had intended to write. After many years working in public health intelligence (PHI), and being increasingly frustrated with my inability to determine if I was making a difference, I began looking for better ways of delivering 'intelligence'. I eventually found 'systems thinking', although not right away, and this book was going to be about my learning journey into systems thinking. The focus was going to be:

\begin{itemize}[leftmargin=1 cm]
	\item My problems, as an analyst, in trying to deliver PHI,
	\item Improving the practice of public health intelligence, and
	\item Produce a user guide for colleagues in my organisation. \end{itemize}

At some point, I embraced this writing project as part of a much wider learning journey that included many colleagues. I've lost count of how many times I've stopped and restarted this writing project. What started as a modest user guide has become a much bigger writing endeavour reflecting the length of our learning journey. This isn't just about me, the analyst, or public health intelligence as a function within public health. I'll explain why later, but this about all of collectively being more effective: Public health intelligence colleagues, public health colleagues, Civil Service colleagues, and especially wider public sector colleagues.

However, before I got to this point there were at least four years of pushing forward and repeatedly running into a wall. Sometimes I had to go back to double check my facts (a lot of that), sometimes I was re-evaluating whether I really believed what I had just written, but mostly it was because what I was writing sounded a bit hollow (maybe readers will still think that). Hollow in the sense that I didn't feel I had any direction to give, even after years of trying to incorporate systems thinking into what I was trying to do.

In retrospect, this project feels like I've been peeling an onion for decades. I was looking for something simple (in all the wrong places), direct and user-guide-ish. So, I got out a knife, and peeled that first layer off. It wasn't looking like a systems thinking user guide at all. I thought I knew what systems thinking was all about, but I hadn't really scratched the surface.

\section{The Progress of Public Health Practice}
I've definitely been an unreasonable man most of my adult life, and my unreasonableness has grown with age and experience. That's not to say that public health hasn't progressed in its thinking, of course it has, but just not by leaps and bounds.

For example, if you are one of my public health colleagues you already know that public health practice is often referred to as: \begin{quote}
	The science and art of promoting and protecting health and wellbeing, preventing ill health, and prolonging life through the organised efforts of society.\sidenote{Faculty of Public Health (fph.org.uk), as of 14 July 2025.}
\end{quote}
As good a definition of our aims as any discipline could have done in so few words. However, what fewer colleagues will know that this quote originates from 1920: 
\begin{quote}
	The science and art of preventing disease, prolonging life, and promoting physical health and efficiency through organized community efforts for the sanitation of the environment, the control of community infections, the education of the individual in principles of personal hygiene, the organization of medical and nursing service for the early diagnosis and preventive treatment of disease, and the development of the social machinery which will ensure to every individual in the community a standard of living adequate for the maintenance of health.\sidenote{Charles-Edward Amory Winslow, an American bacteriologist, 1877-1957.}
\end{quote}
While wordier than our contemporary version, and perhaps without explicit consideration of mental health, it is clear that our aims have remained the same for more than 100 years. Our ambition  remains largely unchanged.

What about our methodology and approach to achieving these aims, have they progressed? I am going to argue that no, we are essentially approaching population the same way we always have. Without a doubt we are now far more technically, scientifically and analytically savvy, and our social views have also evolved, we have been making methodological progress. Nonetheless, our approach overall remains the same, epidemiology, often referred to as the science of public health, still follows the same principles of reductionism, as have most of the sciences.

While we have advanced our sciences to enable deeper and broader understanding of new diseases, forms of ill health, and medical technologies, we no longer appear to be enjoying the same level of achievement vis a vis our aims. Our populations no longer getting healthier each year, progress has slowed and even reversed in some regards. To be sure, our populations are longer lived than they were in 1920, but for how long?

This book is about effectiveness in public health practice and what I, and many colleagues, believe we need to change about our approach if we are to recapture the early progress of our discipline of public health. Furthermore, I am going to suggest how we and our organisations need to change our practice in order to improve effectiveness. You may indeed see this as unreasonable, and you wouldn't be the first of my colleagues to suggest this.

\section{My Personal Journey}
I'm going to relay to you some of my career story and public health experiences, which will probably be different from yours. Nonetheless, I hope it will encourage you to reflect on your career experiences areas of interest, as have I, in order to see them in a new light. I think that this can be a more powerful way of launching your own systems learning journey.

I started questioning my approach to and the value of what I'd been doing in public health intelligence for more than a decade. Data driven organisations. Information for action. Is it just me, or do public health (maybe even public sector) inteventions less effective and less impactfull that they used to be? At the very least, I can't see changes in interventions, decision-making or population health outcomes as a result of the work my colleagues and I deliver.

The questions I've been asking are:
\begin{description}
	\item[Impactfulness] Are our, in the widest public sector sense, interventions having an impact? I honestly don't know, but it feels like we aren't having as much of an impact or as consistent an impact as we might have in the past. I'll take you through what public sector leaders are saying about this though, and spoiler alert, many colleagues believe that rising levels of complexity are at the root of the problem.
	\item[Workforce Effectiveness] Does this meant that we (public health intelligence analysts, the wider public health and public sector workforce) are also not being as effective as we once were? Are we happily lingering in the tepid bath of decline, as one British Prime Minister recently suggested? If this means that we are locking into a single problem solving paradigm of reductionist scientifically based analysis, then yes, we might just have to change how we approach problematic and complex situations in the public sector.
	\item[Systems Thinking] I've heard, and you have too or you wouldn't be here, that there is this thing called systems thinking that could be put into practice that might help us improve our effectiveness and impact within the public sector. It does mean that we have to adopt a new way of working and thinking. This will take some considerable effort on our part, but what other option do we have?
\end{description}

This book is really about trying to answer these and loads more questions that I've been asking myself and responding to some of the many challenges from colleagues. The push-back has comes in various guises, but it seems to have three components:
\begin{enumerate}[leftmargin=1 cm]
	\item Analysis produces the only acceptable evidence-base,
	\item The resulting scientific evidence-base underpins modern public health practice, and
	\item There are no acceptable alternatives to informing practice.
\end{enumerate}
Consequently, some colleagues will argue that we must focus all efforts on improving analysis by adopting more data science methods, including the use of large language models, and moving from analysis to analytics to automate analysis processes, e.g. producing dashboards -- in short, we just need to do our jobs more efficiently in order that we have a bigger impact.

Thus we have a social norm for analysts (maybe most of the public sector workforce) of all stripes, not just public health intelligence analysts, that the reductionist scientific process of analysis produces the only acceptable evidence-based. Challenging this deeply ingrained belief turns out to be a huge challenge in the public sector, maybe even "unreasonable", besides which it will take too long and we don't have enough resources to satisfy the analysis requests.




